\subsection{Distancia de Edición y LCS (DP Bit-Paralelo)}
\label{alg:6-2}

\graybold{Preguntas guía:}

\begin{itemize}
\item ¿Me piden el mínimo de inserciones, borrados y reemplazos para transformar una cadena en otra (distancia de Levenshtein)?
\item ¿Me piden la subsecuencia común más larga (LCS), o la distancia usando solo inserciones y borrados, que es $n + m - 2\cdot\mathrm{LCS}$?
\item ¿La tabla \bigO{nm} tiene decenas de millones de celdas y el DP clásico en Python no entra en el tiempo?
\item ¿Puedo aprovechar que celdas vecinas de la tabla difieren en a lo sumo 1, para codificar una columna entera en bits?
\end{itemize}

\graybold{Utilidad:} Calcula la distancia de edición (y, con una variante, la LCS) entre dos cadenas. El DP clásico llena una tabla de $n \times m$ celdas, lo cual es simple pero inviable en Python cuando $n, m$ llegan a varios miles. La versión bit-paralela representa una columna entera de la tabla como enteros de $m$ bits y la actualiza con unas quince operaciones de bits por carácter; como Python ejecuta esas operaciones sobre enteros grandes en C, el bucle interpretado pasa de $nm$ iteraciones a solo $n$. Es la misma idea que sirve para búsqueda aproximada de patrones (con $k$ errores) y para LCS en tiempo casi lineal.

\graybold{Complejidad:} \bigO{nm} el DP clásico; \bigO{n \cdot \lceil m/w\rceil} el bit-paralelo, con $w = 64$ el tamaño de palabra. En Python esto se traduce en $n$ iteraciones del bucle, cada una con operaciones sobre enteros de $m$ bits resueltas en C.

\graybold{Fundamento teórico:} El DP clásico define $D[i][j]$ como la distancia entre los prefijos $a[0..i)$ y $b[0..j)$, con $D[i][0] = i$, $D[0][j] = j$ y $D[i][j] = \min(D[i-1][j] + 1,\ D[i][j-1] + 1,\ D[i-1][j-1] + [a_{i-1} \ne b_{j-1}])$: borrar, insertar, o reemplazar (gratis si las letras coinciden). La observación que habilita el paralelismo es que dos celdas vecinas de la tabla, horizontal o verticalmente, difieren en EXACTAMENTE $-1$, $0$ o $+1$: agregar un carácter a una de las cadenas nunca cambia la distancia en más de una operación. Por lo tanto, en lugar de guardar los valores de una columna, alcanza con guardar sus DIFERENCIAS verticales, codificadas en dos máscaras de bits: \texttt{Pv} marca las filas donde la diferencia es $+1$ y \texttt{Mv} las filas donde es $-1$ (las demás valen 0). El algoritmo de Myers, adaptado por Hyyrö a la distancia global, deriva de la recurrencia unas pocas fórmulas de bits que calculan, a partir de la columna anterior y de la máscara \texttt{Eq} de posiciones de $a$ donde aparece la letra actual de $b$, las diferencias horizontales (\texttt{Ph}, \texttt{Mh}) y luego la nueva columna. El paso no evidente es la SUMA \texttt{(Eq \& Pv) + Pv}: el acarreo de la suma binaria se propaga a lo largo de tramos consecutivos de bits, y eso modela exactamente cómo una coincidencia en la diagonal se propaga hacia abajo por la columna, algo que celda por celda requeriría un bucle. El valor final no se reconstruye entero: se lleva solo \texttt{score}, el valor de la última fila, que empieza en $m$ y en cada columna sube o baja según el bit más alto de \texttt{Ph} o \texttt{Mh}. El desplazamiento \texttt{(Ph <{}< 1) | 1} introduce un $+1$ en la fila 0 porque $D[0][j] = j$ crece en uno por columna; esa es la diferencia con la búsqueda aproximada de patrones, donde se introduce un 0 para permitir que la coincidencia empiece en cualquier lugar del texto. Para la LCS la situación es aún más simple: las diferencias verticales son solo 0 o 1, así que basta una máscara, y la actualización se reduce a una suma y una resta. En Python hay un detalle adicional: el operador \texttt{\textasciitilde} sobre enteros produce negativos (es $-x-1$), por lo que toda negación y todo desplazamiento se recortan con \texttt{\& mask} a $m$ bits.~\cite{myers1999}

\graybold{Ejercicio aplicado}

(CSES 1639 — Edit Distance) Dadas dos cadenas en mayúsculas de longitudes $n, m \le 5000$, calcular su distancia de edición con inserciones, borrados y reemplazos.

\graybold{I/O:} Dos líneas, una cadena cada una $\Rightarrow$ \texttt{sys.stdin.buffer.readline} (patrón 1), trabajando sobre \texttt{bytes} para que cada letra sea un entero usable directamente como clave de la tabla \texttt{peq}. Sobre el rendimiento conviene ser explícito, porque es la razón de ser de esta versión: el DP clásico, incluso con filas rodadas y sin llamadas a \texttt{min}, tarda \aprox{}3.3s en el caso $5000 \times 5000$, muy por encima del límite de 1 segundo; la versión bit-paralela resuelve el mismo caso en \aprox{}0.03s, unas 125 veces más rápido, con la misma respuesta. Verificado contra el caso de ejemplo y contrastado con el DP clásico en 6069 casos, incluyendo alfabetos de 1 a 26 letras, cadenas iguales, prefijos, cadenas sin letras en común, el clásico \texttt{KITTEN}/\texttt{SITTING}, y 60 pares de entre 65 y 300 caracteres para ejercitar el acarreo entre palabras de 64 bits, sin fallos.

\graybold{Código solución}

\begin{lstlisting}[language=Python]
import sys

def edit_distance(a, b):
    m = len(a)
    mask = (1 << m) - 1         # recorta a m bits: ~x en Python da negativos
    alto = 1 << (m - 1)         # bit de la ultima fila de la tabla
    # peq[c] = mascara de las posiciones de a donde aparece la letra c
    peq = {}
    for i, c in enumerate(a):
        peq[c] = peq.get(c, 0) | (1 << i)
    Pv = mask                   # columna 0: D[i][0] = i, todas las diferencias son +1
    Mv = 0
    score = m                   # valor de la ultima fila, D[m][0] = m
    for c in b:
        Eq = peq.get(c, 0)
        Xv = Eq | Mv
        # la SUMA propaga las coincidencias de la diagonal hacia abajo por acarreo
        Xh = (((Eq & Pv) + Pv) ^ Pv) | Eq
        Ph = Mv | (~(Xh | Pv) & mask)   # diferencias horizontales +1
        Mh = Pv & Xh                    # diferencias horizontales -1
        if Ph & alto:
            score += 1
        elif Mh & alto:
            score -= 1
        # el |1 mete el +1 de la fila 0 (D[0][j] = j): distancia GLOBAL
        Ph = ((Ph << 1) | 1) & mask
        Mh = (Mh << 1) & mask
        Pv = Mh | (~(Xv | Ph) & mask)   # nueva columna de diferencias verticales
        Mv = Ph & Xv
    return score

def solve():
    a = sys.stdin.buffer.readline().strip()
    b = sys.stdin.buffer.readline().strip()
    print(edit_distance(a, b))

solve()
\end{lstlisting}

\graybold{Extras: LCS bit-paralela} — la subsecuencia común más larga con la misma técnica. Como las diferencias verticales de la tabla de LCS son solo 0 o 1, basta una máscara y la actualización se reduce a una suma y una resta. Verificada contra una tabla de referencia en 5050 casos; en $5000 \times 5000$ tarda \aprox{}0.004s.

\begin{lstlisting}[language=Python]
def lcs_bits(a, b):
    m = len(a)
    mask = (1 << m) - 1
    peq = {}
    for i, c in enumerate(a):
        peq[c] = peq.get(c, 0) | (1 << i)
    V = mask
    for c in b:
        U = V & peq.get(c, 0)
        V = ((V + U) | (V - U)) & mask
    return m - bin(V).count("1")    # ceros de V = largo de la LCS
                                    # (en Python 3.10+: V.bit_count())

# con solo insertar y borrar (sin reemplazo): distancia = n + m - 2 * lcs_bits(a, b)
\end{lstlisting}
